% !PRINTED 464
\weberpage{464}
%<<CONT
proved, the degree $n_1$ of the field $\Omega(\beta)$ is a
divisor of $n$. If now $\Omega'$ is not exhausted by $\Omega(\beta)$,
we take a quantity $\gamma$ in $\Omega'$, but not in
$\Omega(\beta)$; then the field $\Omega(\beta,
\gamma)$ is a divisor of $\Omega'$ and of $\Omega(\alpha)$ and has,
for its part, $\Omega(\beta)$ as a divisor. The degree of $\Omega(\beta,
\gamma)$ is therefore greater than that of $\Omega(\beta)$ and smaller
than that of $\Omega(\alpha)$. If $\Omega'$ is not yet exhausted
by this, we continue in the same way, but must finally come to an
end, since the degrees of the fields $\Omega(\beta), \Omega(\beta,
\gamma) \ldots$ always increase and yet remain smaller than $n$.
From this it follows:

\gesperrt{5. Every divisor of $\Omega(\alpha)$ which contains the
field $\Omega$ is an algebraic field $\Omega(\beta)$ over
$\Omega$. The degree of $\Omega(\beta)$ is a divisor of the degree
of $\Omega(\alpha)$, and if both fields have the same degree,
they are identical.}

We can therefore also state our definition thus:

\gesperrt{6. An algebraic field over $\Omega$ is primitive
if it has no field over $\Omega$ except $\Omega$ and itself
as a divisor.}

\begin{center}
\S.\,145.

\gesperrt{Normal fields. Galois resolvent.}
\end{center}

If $\Omega(\alpha)$ is an algebraic field, then the
conjugate fields
% [sic] The source prints alpha_{m-1}; §144 defines the conjugate fields only through alpha_{n-1}, where n is the irreducible-factor degree. Preserved diplomatically.
\[
\Omega(\alpha), \Omega(\alpha_1) \ldots \Omega(\alpha_{m-1})
\]
are all of the same degree, and hence if one of them is contained
in another, the two are identical.

A field which is identical with all its conjugate fields
is called a \gesperrt{normal field}. In the normal fields
much simpler laws prevail, and the great advance which
algebra owes to \gesperrt{Galois} rests essentially
on the fact that arbitrary fields are reduced to normal
fields. The normal fields are therefore also called \gesperrt{Galois
fields}.

If a field $\Omega(\varrho)$ of degree $\mu$ has the
property that the numbers $\varrho_1,
\varrho_2 \ldots \varrho_{\mu-1}$ conjugate to $\varrho$ are all
contained in $\Omega(\varrho)$, then it is a normal field; for then the
fields $\Omega(\varrho_1), \Omega(\varrho_2) \ldots
\Omega(\varrho_{\mu-1})$ are also all contained in $\Omega(\varrho)$
and hence all identical with $\Omega(\varrho)$.
